% Avsnitt 14.8, ur lectures/L14/appendix/b_exercises.md.

\section{Övningar}
\label{c14:sec:exercises}

\begin{aside}
De här övningarna befäster kapitlet. Konstruera din egen \file{tx_shift_reg.vhd} utifrån
specifikationen i \secref{c14:sec:iface} till \ref{c14:sec:trace}; den utdelade testbänken, körd
enligt bilaga~\ref{app:sim}, är kontrollen av att den beter sig korrekt.
\end{aside}

\exercise{\code{tx_shift_reg}}{VHDL}
\label{c14:ex:txsr}
Skriv din egen \file{tx_shift_reg.vhd} utifrån kapitlet. Verifiera den:

\begin{shell}
cd controller
ghdl -a --std=93 can_def.vhd tx_shift_reg.vhd tx_shift_reg_tb.vhd
ghdl -e --std=93 tx_shift_reg_tb
ghdl -r --std=93 tx_shift_reg_tb --assert-level=error
\end{shell}

Om en kontroll fallerar, läs assertion-meddelandet; det namnger exakt vilket skift det gäller och
vad som förväntades.

\exercise{Följ en stoppad sekvens för hand, sedan i simulering}{Simulering}
\label{c14:ex:trace}
\xpart{a} Följ för hand \code{tx_shift_reg} laddad med \code{data = "11100000"},
\code{bit_count = 8}, med start vid en gruppgräns där den föregående gruppens sista bit var \code{0}
med en följdlängd på 3 (alltså \code{consecutive = 3} och \code{last_bit = '0'} in i gruppen). Skriv
ut, bit för bit: vad som presenteras vid \code{load}, och sedan vid varje efterföljande
\code{shift}, med varje instoppad stoppbit markerad och med angivande av vilket skift som höjer
\code{done}.

\xpart{b} Bekräfta din handföljning i simulering: skriv en kort testbänk (eller återanvänd strukturen
i \file{tx_shift_reg_tb.vhd}) som förladdar samma startvillkor med följdlängd 3, sedan laddar
\code{"11100000"} och vid varje skift \code{assert}:ar \code{tx_bit}, \code{stuff} och \code{done}
mot värdena du förutsade (eller \code{report}:ar dem så att du kan läsa dem). Stämmer dess utfall av
passera eller fallera med din handföljning?

\exercise{Enbitsgruppen som ändå är skyldig en stoppbit}{Simulering}
\label{c14:ex:onebit}
\Secref{c14:sec:timing}:s sista punkt namnger ett gränsfall utan att arbeta igenom det: \emph{om
\code{bit_count = 1} och den enda biten utlöser en stoppbit, väntar \code{done} ändå på att
stoppbiten skickas först.} Det är inte hypotetiskt; \code{can_controller} laddar enbitsgrupper (SOF)
och fyrabitsgrupper (ID:ts undre halva plus RTR), så en grupp kan verkligen sluta mitt i en följd.

Sätt upp det precis som test 2 och 3 i \code{tx_shift_reg_tb} gör, utan reset emellan: ladda
\code{"00001111"} med \code{bit_count = 8} och skifta ut den helt, så att följdräknaren står på fyra
\code{'1'}:or i rad. Ladda sedan en \textbf{enbitsgrupp} vars enda bit också är \code{'1'}
(\code{data = "1-------"}, \code{bit_count = 1}).

\xpart{a} Följ vad som händer, presentation för presentation: vad \code{load} lägger på
\code{tx_bit}, vad följdräknaren når, och sedan vad varje följande \code{shift} gör. Markera vilken
presentation som bär stoppbiten och vilken som höjer \code{done}.

\xpart{b} Hur många \code{shift}-pulser förbrukar den här enbitsgruppen totalt, och hur många skulle
det ha förbrukat om följdräknaren i stället stått på noll? Säg var skillnaden kommer ifrån.

\xpart{c} Anta att \code{done} hade fått pulsa så snart gruppens enda riktiga bit presenterats, utan
hänsyn till den väntande stoppbiten. \code{can_controller} skulle då ladda nästa grupp under nästa
bitperiod. Beskriv exakt vad den mottagande noden i så fall skulle se på ledningen, och vilken av
\code{rx_shift_reg}:s utgångar som till slut skulle rapportera det (\chapref{c15:ch}).

\xpart{d} Bekräfta din genomgång i simulering, genom att utöka mönstret från test 2 och test 3 i
\file{tx_shift_reg_tb.vhd} med den här enbitsgruppen som ett fjärde fall.

\exercise{Varifrån \code{track_run} anropas}{Simulering}
\label{c14:ex:trackrun}
\Secref{c14:sub:tracker} lägger följdräknaren i en procedur och är tydlig med var den proceduren får
anropas: från en gren som placerar en bit, aldrig från processens toppnivå. Båda halvorna av den
regeln är värda att bryta med flit, i en kladdkopia av din modul.

\xpart{a} Flytta ut anropet \code{track_run(data(7));} ur grenen \code{load}, så att \code{load}
låser gruppen och presenterar dess första bit men räknaren aldrig ser den biten. Förutsäg vad
följdräknaren då visar under resten av ramen, och vilket av testbänkens tre tester som fallerar
först. Kör det.

\xpart{b} Lägg nu ett enda \code{track_run(tx_bit);} högst upp i grenen
\code{elsif (rising_edge(clock))}, ovanför standardvärdena, och ta bort alla tre anropen inuti
grenarna. \code{tx_bit} är en \code{out}-port, så börja med att säga varför det inte ens kompilerar,
och upprepa sedan experimentet med en intern kopia av den sända biten i dess ställe. Med
\code{TICKS_PER_BIT} satt till 50, hur många gånger körs räknaren nu per CAN-bit, och vad gör
\code{consecutive} på den andra klockflanken i den allra första bitperioden?

\xpart{c} \Secref{c14:sub:tracker} säger också att proceduren måste deklareras \emph{inuti} processen
i stället för i arkitekturen. Flytta ut den till arkitekturens deklarativa del, oförändrad, och
analysera filen. Den kompilerar inte; läs felen, ett per signaltilldelning i proceduren, innan du går
vidare. Gör sedan en version på arkitektursnivå som \emph{faktiskt} kompilerar, med identiskt
beteende och en testbänk som passerar. Jämför de två anropsställena och säg vad den andra versionen
kostar, och vilken ny fråga den väcker som versionen inuti processen inte kan väcka.
